\section{Methods: Group 6 --- Assessement}
\label{sec:g6}

The assessement layer turns raw measurements into status, severity, and flags by comparing each value
to a central, fully cited threshold catalog. It is deliberately separated from measurement extraction so
that the science (the thresholds) lives in one auditable place.

\paragraph{The threshold catalog.} \texttt{thresholds.py} holds one specification per measurement key
(clinical name, unit, severity bands, flag threshold, citation, modality caveat, and a tag of
raw/derived/quality). \texttt{classify(key, value)} walks the ordered bands and returns a standardized
status (\texttt{within\_reference} / \texttt{outside\_reference} / \texttt{review\_only} /
\texttt{not\_interpretable}), a per-measurement severity label, the boolean clinical flag
(\texttt{true} iff the value lands in an ``outside'' band), and the locked citation and caveat. Where no
cited threshold exists, the spec says so (citation \texttt{None}, status \texttt{review\_only}) --- the
system never invents a number.

\paragraph{Provenance and honesty rules.} Every clinical number carries a primary-source citation
(PMID/DOI), copied verbatim from verified-research notes because automated literature search can
hallucinate author names on real identifiers. Modality caveats are attached where a threshold is a
radiograph/CT/lumbar borrow (e.g.\ Torg, the osseous canal cuts, the lumbar-origin disc rules), so the
report never silently conflates a bony-canal threshold with a soft-tissue MRI measurement.

\paragraph{Quality versus clinical flags.} A flag is classified as a quality/caution flag (not a patient
abnormality) if its name matches markers such as \texttt{outlier}, \texttt{unreliable},
\texttt{approximate}, \texttt{low\_confidence}, \texttt{misaligned}, \texttt{resolution}, or
\texttt{warning}. This fixes an earlier mis-classification where a vertebral-tilt quality flag could be
read as pathology.

\paragraph{Syndrome indicators.} Provisional, advisory-only combination rules sketch myelopathy (canal
narrowing $+$ low SAC $+$ cord signal at a level) and radiculopathy (weak on sagittal-only MRI, since
foraminal dimensions are not measured) --- each marked provisional and never diagnostic.

\paragraph{Threshold corrections folded in from validation research.} A dedicated research pass produced
corrections now reflected as caveats/bands: the dural-sac/SAC/Torg cuts are osseous/radiograph borrows
that over-flag on MRI; the \SI{1.35}{\milli\meter} disc ``cord-risk'' figure is unverified; the Miyazaki
grade-IV is non-discriminating in isolation; and a prior MRI-vs-CT disc-height offset had been inverted.

\paragraph{Status.} The engine and catalog are built and unit-tested; the end-to-end run is gated on the
measurement validation that this paper reports, after which the gap-filling thresholds (MSCC/MCC,
per-level segmental norms, etc.) are the remaining Group~6 work.
